% ═══════════════════════════════════════════════════════════════
% ARBEITSBLATT - Physik Klasse 6: Verkehrssicherheit
% 27 BE (★) | 30 BE (★★) | 33 BE (★★★) | 2 Seiten
% ═══════════════════════════════════════════════════════════════

\documentclass[11pt, a4paper]{article}

\usepackage[utf8]{inputenc}
\usepackage[T1]{fontenc}
\usepackage[ngerman]{babel}
\usepackage{helvet}
\renewcommand{\familydefault}{\sfdefault}

\usepackage[margin=1.5cm, top=1.8cm, bottom=1.2cm]{geometry}
\usepackage{enumitem}
\usepackage{tabularx}
\usepackage{array}
\usepackage{colortbl}
\usepackage{tikz}
\usepackage{amssymb}
\usepackage{amsmath}
\usepackage{fancyhdr}
\usepackage{xcolor}
\usepackage{tcolorbox}
\tcbuselibrary{skins, breakable}
\usepackage{qrcode}
\usepackage[hidelinks]{hyperref}

% Lernpfad-URL
\newcommand{\lpurl}{https://devbox42.github.io/sciencesim/physik/kl06/02-bewegung/std-04-verkehrssicherheit/LP-04-verkehrssicherheit.html}

% Farben
\definecolor{physik}{HTML}{2c5aa0}
\definecolor{physik-hell}{HTML}{e8f0fa}
\definecolor{hellgrau}{HTML}{F5F5F5}
\definecolor{orange-warn}{HTML}{b35c00}
\definecolor{orange-hell}{HTML}{fff3e0}
\definecolor{reaktion}{HTML}{2196f3}
\definecolor{brems}{HTML}{ff9800}
\definecolor{anhalt}{HTML}{4caf50}

% Kopfzeile
\pagestyle{fancy}
\fancyhf{}
\fancyhead[L]{\small\textcolor{physik}{\textbf{Physik Klasse 6}} | Arbeitsblatt: Verkehrssicherheit}
\fancyhead[R]{\small Name: \underline{\hspace{3.5cm}}}
\fancyfoot[R]{\small\thepage/2}
\renewcommand{\headrulewidth}{0pt}

% Kompakte Boxen
\newtcolorbox{aufgabenbox}[1][]{
    colback=white,
    colframe=physik,
    boxrule=0.8pt,
    arc=0pt,
    left=6pt, right=6pt, top=4pt, bottom=4pt,
    title={\small\textbf{#1}},
    coltitle=black,
    fonttitle=\sffamily,
    colbacktitle=white,
    boxed title style={arc=0pt, boxrule=0.8pt, colframe=physik}
}

\newtcolorbox{infobox}{
    colback=orange-hell,
    colframe=orange-warn,
    boxrule=1pt,
    arc=0pt,
    left=6pt, right=6pt, top=6pt, bottom=4pt
}

% Befehle
\newcommand{\punkte}[1]{\hfill\textcolor{gray}{\small /#1}}
\newcommand{\luecke}[1]{\underline{\hspace{#1}}}
\newcommand{\sternchen}[1]{\textcolor{physik}{\textbf{(#1)}}}

% ISO 80000
\newcommand{\fz}[1]{\textit{\textrm{#1}}}
\newcommand{\einheit}[1]{\textsf{\textup{#1}}}
\newcommand{\mps}{$\tfrac{\textsf{m}}{\textsf{s}}$}

\setlength{\parindent}{0pt}
\setlength{\parskip}{2pt}

\begin{document}

% ═══════════════════════════════════════════════════════════════
% KOPFBEREICH (kompakt)
% ═══════════════════════════════════════════════════════════════

\noindent
\begin{minipage}[t]{0.78\textwidth}
\textcolor{physik}{\rule{\textwidth}{2pt}}\\[0.2em]
{\Large\bfseries\textcolor{physik}{Verkehrssicherheit: Anhalteweg}}\\[0.1em]
\textcolor{physik}{\rule{\textwidth}{0.5pt}}
\end{minipage}%
\hfill
\begin{minipage}[t]{0.18\textwidth}
\raggedleft
\href{\lpurl}{\qrcode[height=1.5cm]{\lpurl}}\\[-2pt]
{\tiny\textcolor{gray}{Lernpfad}}
\end{minipage}

\vspace{-0.5em}

\begin{flushright}
\small
\begin{tabular}{|c|c|c|c|}
\hline
\cellcolor{physik-hell} \textbf{$\star$} & \cellcolor{physik-hell} \textbf{$\star\star$} & \cellcolor{physik-hell} \textbf{$\star\star\star$} & \cellcolor{physik-hell} \textbf{Erreicht} \\
\hline
/27 BE & /30 BE & /33 BE & \hspace{1cm} \\
\hline
\end{tabular}
\end{flushright}

\vspace{-0.3em}

% ═══════════════════════════════════════════════════════════════
% BREMSWEG-TABELLE (kompakt)
% ═══════════════════════════════════════════════════════════════

\begin{infobox}
\textbf{\textcolor{orange-warn}{Bremsweg-Tabelle}}

\vspace{0.2em}

\renewcommand{\arraystretch}{1.2}
\small
\begin{tabularx}{\textwidth}{|>{\centering\arraybackslash}p{2.8cm}|>{\centering\arraybackslash}X|>{\centering\arraybackslash}X|>{\centering\arraybackslash}X|}
\hline
\rowcolor{physik-hell}
\textbf{Geschwindigkeit} & \textbf{Trocken} & \textbf{Nass} & \textbf{Eis} \\
\hline
5 \mps{} (18 km/h) & 2 m & 4 m & 10 m \\
\hline
10 \mps{} (36 km/h) & 5 m & 10 m & 25 m \\
\hline
15 \mps{} (54 km/h) & 10 m & 20 m & 50 m \\
\hline
20 \mps{} (72 km/h) & 20 m & 40 m & 100 m \\
\hline
30 \mps{} (108 km/h) & 45 m & 90 m & 225 m \\
\hline
\end{tabularx}
\end{infobox}

\vspace{0.2em}

% ═══════════════════════════════════════════════════════════════
% MEINE REAKTIONSZEIT (Mini-Experiment)
% ═══════════════════════════════════════════════════════════════

\noindent
\begin{minipage}[t]{0.48\textwidth}
\small\textbf{Meine Reaktionszeit} (aus dem Lernpfad):\\[0.2em]
\hspace{0.5cm} $t_R$ = \luecke{2cm} s
\end{minipage}%
\hfill
\begin{minipage}[t]{0.48\textwidth}
\small\textit{Typisch: 0,2--0,5 s (Ampel-Test)\\
Beim Autofahren rechnet man mit 1 s.}
\end{minipage}

\vspace{0.3em}

% ═══════════════════════════════════════════════════════════════
% SKIZZE (neu)
% ═══════════════════════════════════════════════════════════════

\begin{aufgabenbox}[Skizze: Anhalteweg \sternchen{$\star$} \punkte{3}]
\textbf{Ergänze} die Skizze: Beschrifte die Strecken und zeichne die Balken passend ein.

\vspace{0.3em}

\begin{center}
\begin{tikzpicture}[scale=0.85]
    % Rahmen
    \draw[thick] (0,0) rectangle (14,2.2);

    % Auto (einfach)
    \fill[physik] (0.3,0.8) rectangle (1.1,1.3);
    \fill[physik] (0.4,1.3) rectangle (0.95,1.6);
    \fill[black] (0.45,0.75) circle (0.12);
    \fill[black] (0.95,0.75) circle (0.12);

    % Gefahr-Symbol
    \fill[red] (13,1.1) circle (0.25);
    \node at (13,1.1) {\tiny\textbf{!}};

    % Balken-Vorlage (gestrichelt)
    \draw[dashed, gray] (1.3,1.0) -- (6,1.0);
    \draw[dashed, gray] (6,1.0) -- (12,1.0);

    % Beschriftungslinien
    \draw[<->, thick, reaktion] (1.3,0.4) -- (6,0.4);
    \draw[<->, thick, brems] (6,0.4) -- (12,0.4);
    \draw[<->, thick, anhalt] (1.3,1.8) -- (12,1.8);

    % Lücken für Beschriftung
    \node at (3.65,0.15) {\small\luecke{2.5cm}};
    \node at (9,0.15) {\small\luecke{2.5cm}};
    \node at (6.65,2.0) {\small\luecke{2.5cm}};

    % Formel rechts
    \node[anchor=west] at (12.3,1.1) {\small $s_A = s_R + s_B$};
\end{tikzpicture}
\end{center}
\end{aufgabenbox}

\vspace{0.3em}

% ═══════════════════════════════════════════════════════════════
% AUFGABE 1: Begriffe (4 BE) - kompakt
% ═══════════════════════════════════════════════════════════════

\begin{aufgabenbox}[Aufgabe 1: Begriffe \sternchen{$\star$} \punkte{4}]
\textbf{Ergänze} mit: \textit{Schrecksekunde, Reaktionsweg, Bremsweg, Anhalteweg}

\vspace{0.3em}

\begin{enumerate}[label=\alph*), itemsep=0.4em, leftmargin=*]
\item Die Zeit vom Erkennen einer Gefahr bis zum Beginn der Reaktion heißt \luecke{3cm}.
\item Der \luecke{2.5cm} ist die Strecke während der Reaktionszeit.
\item Vom Beginn der Bremsung bis zum Stillstand: \luecke{2.5cm}.
\item Die \textbf{gesamte} Strecke vom Erkennen bis zum Stillstand ist der \luecke{2.5cm}.
\end{enumerate}
\end{aufgabenbox}

\vspace{0.3em}

% ═══════════════════════════════════════════════════════════════
% AUFGABE 2: Reaktionsweg (6 BE) - kompakt
% ═══════════════════════════════════════════════════════════════

\begin{aufgabenbox}[Aufgabe 2: Reaktionsweg berechnen \sternchen{$\star$} \punkte{6}]
\textbf{Berechne} den Reaktionsweg. Reaktionszeit: \fz{t} = 1 \einheit{s}. \hfill \textit{Formel: $s_R = v \cdot t$}

\vspace{0.3em}

\renewcommand{\arraystretch}{1.6}
\begin{tabularx}{\textwidth}{|X|X|X|}
\hline
\textbf{a)} Fahrrad, \fz{v} = 5 \mps{} & \textbf{b)} Auto, \fz{v} = 10 \mps{} & \textbf{c)} Auto, \fz{v} = 20 \mps{} \\
\hline
$s_R$ = & $s_R$ = & $s_R$ = \\[0.8em]
\hline
\end{tabularx}
\end{aufgabenbox}

\vspace{0.3em}

% ═══════════════════════════════════════════════════════════════
% AUFGABE 3: Bremsweg ablesen (4 BE) - kompakt
% ═══════════════════════════════════════════════════════════════

\begin{aufgabenbox}[Aufgabe 3: Bremsweg ablesen \sternchen{$\star$} \punkte{4}]
\textbf{Lies} den Bremsweg aus der Tabelle ab.

\vspace{0.3em}

\renewcommand{\arraystretch}{1.5}
\small
\begin{tabularx}{\textwidth}{|X|c|>{\raggedleft\arraybackslash}p{2.2cm}||X|c|>{\raggedleft\arraybackslash}p{2.2cm}|}
\hline
\rowcolor{physik-hell}
\textbf{Situation} & \textbf{v} & $s_B$ & \textbf{Situation} & \textbf{v} & $s_B$ \\
\hline
a) trocken & 10 \mps{} & \einheit{m} & c) trocken & 20 \mps{} & \einheit{m} \\
\hline
b) nass & 10 \mps{} & \einheit{m} & d) Eis & 20 \mps{} & \einheit{m} \\
\hline
\end{tabularx}
\end{aufgabenbox}

\vspace{0.3em}

% ═══════════════════════════════════════════════════════════════
% AUFGABE 4: Anhalteweg (8 BE) - kompakt
% ═══════════════════════════════════════════════════════════════

\begin{aufgabenbox}[Aufgabe 4: Anhalteweg berechnen \sternchen{$\star$} \punkte{10}]
\textbf{Berechne} den Anhalteweg. \hfill \textit{Formel: Anhalteweg = Reaktionsweg + Bremsweg}

\vspace{0.4em}

\textbf{a)} \textbf{Fahrrad}, \fz{v} = 10 \mps{}, \textbf{trockene} Straße \hfill (7 BE)

\vspace{0.2em}
\hspace{0.5cm} $s_R$ = \luecke{1cm} $\cdot$ \luecke{0.8cm} = \luecke{1.2cm} \hfill $s_B$ (Tab.) = \luecke{1.2cm} \hfill $s_A$ = \luecke{1cm} + \luecke{1cm} = \luecke{1.2cm}

\vspace{0.6em}

\textbf{b)} \textbf{Auto}, \fz{v} = 20 \mps{}, \textbf{nasse} Straße \hfill (3 BE)

\vspace{0.2em}
\hspace{0.5cm} $s_R$ = \luecke{2cm} \hfill $s_B$ (Tab.) = \luecke{1.5cm} \hfill $s_A$ = \luecke{2.5cm}
\end{aufgabenbox}

\vspace{0.3em}

% ═══════════════════════════════════════════════════════════════
% AUFGABE 5: Vergleich (3 BE) - ★★
% ═══════════════════════════════════════════════════════════════

\begin{aufgabenbox}[Aufgabe 5: Vergleich \sternchen{$\star\star$} \punkte{3}]
Emma fährt mit \fz{v} = 10 \mps{}, Felix mit \fz{v} = 20 \mps{}. Beide bremsen auf \textbf{trockener} Straße.

\vspace{0.4em}

\textbf{a)} Um wie viel Meter ist der Anhalteweg von Felix \textbf{länger}? \hfill (2 BE)

\vspace{0.3em}

\renewcommand{\arraystretch}{1.4}
\begin{tabularx}{\textwidth}{|c|c|>{\centering\arraybackslash}p{2.5cm}|c|>{\centering\arraybackslash}p{2.5cm}|}
\hline
\rowcolor{physik-hell}
& $v$ & $s_R$ & $s_B$ (Tab.) & $s_A$ \\
\hline
Emma & 10 \mps{} & & & \\
\hline
Felix & 20 \mps{} & & & \\
\hline
\end{tabularx}

\vspace{0.3em}

\hspace{0.5cm} Differenz: $s_A$(Felix) $-$ $s_A$(Emma) = \luecke{3cm}

\vspace{0.3em}

\textbf{b)} Was fällt dir auf? \hfill (1 BE)

\vspace{0.6cm}
\end{aufgabenbox}

\vspace{0.3em}

% ═══════════════════════════════════════════════════════════════
% AUFGABE 6: Bewertung (3 BE) - ★★★
% ═══════════════════════════════════════════════════════════════

\begin{aufgabenbox}[Aufgabe 6: Verkehrssicherheit bewerten \sternchen{$\star\star\star$} \punkte{3}]
Es hat gerade angefangen zu regnen. Du fährst mit dem Fahrrad (\fz{v} = 10 \mps{}).

\vspace{0.4em}

\textbf{a)} Berechne den Anhalteweg bei \textbf{nasser} Straße. \hfill (1 BE)

\vspace{0.3em}
\hspace{0.5cm} $s_A$ = \luecke{5cm}

\vspace{0.5em}

\textbf{b)} Was solltest du tun, um sicher unterwegs zu sein? Nenne \textbf{eine} Maßnahme und \textbf{begründe}! \hfill (2 BE)

\vspace{0.3em}
\hspace{0.5cm} Maßnahme: \luecke{10cm}

\vspace{0.4cm}
\hspace{0.5cm} Begründung:

\vspace{0.8cm}
\end{aufgabenbox}

\vfill

\begin{center}
\small\textcolor{gray}{Physik Klasse 6 | Verkehrssicherheit | AB-04 | Seite 2/2}
\end{center}

\end{document}
